\clearpage
\VerdierPhysicalPageAnchor{0230}%
\begingroup\small\setlength{\parskip}{0pt}\setlength{\emergencystretch}{1.5em}
\setlength{\abovedisplayskip}{4pt plus 2pt minus 2pt}%
\setlength{\belowdisplayskip}{4pt plus 2pt minus 2pt}%

\VerdierSemanticAnchor{ch4d-4-1-6-cont-2}
\noindent 次に、

\VerdierSemanticAnchor{ch4d-4-1-6-6}
\begin{equation}
\int_n f\bigl((p,q),(p',q')\bigr)
=h\bigl((p,q),(p',q')\bigr).
\tag{4.1.6.6}
\end{equation}

と置く。さらに、

\VerdierSemanticAnchor{ch4d-4-1-6-7}
\begin{equation}
{}^{m}\!d(p,q,r)\colon{}^{m}\!Y(q,r)
\longrightarrow{}^{m}\!Y(p,q)[1]
\tag{4.1.6.7}
\end{equation}

を合成射

\VerdierSemanticAnchor{ch4d-4-1-6-8}
\begin{equation}
\begin{aligned}
{}^{m}\!Y(q,r)
&\xrightarrow{\int_n g(p,q,r)}
\int_n T_m\bigl({}^{m}\!X(p,q)\bigr)
\xrightarrow{(*)}
\left(\int_n{}^{m}\!X(p,q)\right)[1]\\
&=Y(p,q)[1].
\end{aligned}
\tag{4.1.6.8}
\end{equation}

として定める。ここで射 $(*)$ は、付随する単純複体関手とシフトとの可換
同型である（第 I 章、2.3.2.1）。

こうして単純複体 $Y=\int_n X$ のフィルトレーション

\VerdierSemanticAnchor{ch4d-4-1-6-9}
\begin{equation}
0={}^{m}\!Y(-\infty)\subset\cdots\subset{}^{m}\!Y(q)
\subset\cdots\subset{}^{m}\!Y(+\infty)=Y
\tag{4.1.6.9}
\end{equation}

が得られる。これを\emph{$X$ の第 $m$ 次数によって定まるフィルトレーション}
という。

(4.1.1) で思い出した方法により、フィルター付き複体 (4.1.6.9) に、
$K(\mathcal A)$ に値をとるスペクトル対象

\VerdierSemanticAnchor{ch4d-4-1-6-10}
\begin{equation}
\bigl(\mathcal Y(p,q),\delta(p,q,r)\bigr).
\tag{4.1.6.10}
\end{equation}

を対応させることができる。最後に、次の命題では記法を濫用して、
$\operatorname{comp}(\mathcal A)$ の射とその $K(\mathcal A)$ における像を
同じ記法で表す。

\VerdierSemanticAnchor{ch4d-4-1-7}
\medskip\noindent\textbf{命題 4.1.7.}\quad

\VerdierSemanticAnchor{ch4d-4-1-7-a}
\noindent\textit{a) 対象 ${}^{m}\!Y(p,q)$ (4.1.6.5)、射
$h\bigl((p,q),(p',q')\bigr)$ (4.1.6.6)、および射
${}^{m}\!d(p,q,r)$ (4.1.6.7) は、$K(\mathcal A)$ に値をとり、
$n$-重複複体 $X$ に関手的に依存するスペクトル対象を定める。}

\noindent\textit{このスペクトル対象を、第 $m$ 次数でフィルター付けされた複体
$X$ によって定まる、$K(\mathcal A)$ に値をとるスペクトル対象という。}

\VerdierSemanticAnchor{ch4d-4-1-7-b}
\noindent\textit{b) スペクトル対象
$\bigl(\mathcal Y(p,q),\delta(p,q,r)\bigr)$ (4.1.6.10) と
$\bigl({}^{m}\!Y(p,q),{}^{m}\!d(p,q,r)\bigr)$ の間には標準同型が存在する
（すなわち、$\operatorname{comp}^n(\mathcal A)$ の対象 $X$ に関して
関手的な同型である）。}
\endgroup
